%==============================================================================
% Chapitre 31 - Rotation et stabilisation gyroscopique
%==============================================================================

\section{Introduction}

Les projectiles modernes d'armes à feu possèdent généralement une forme
allongée.

Cette géométrie offre d'excellentes performances aérodynamiques mais ne
garantit pas à elle seule une stabilité suffisante.

Pour maintenir une orientation correcte pendant le vol, les projectiles sont
mis en rotation rapide autour de leur axe longitudinal.

Cette technique est appelée stabilisation gyroscopique.

\begin{aretenir}

La rotation d'un projectile constitue le principal mécanisme de stabilité des
balles modernes.

\end{aretenir}

\section{Origine de la rotation}

La rotation est produite par les rayures du canon.

Ces rayures imposent au projectile un mouvement hélicoïdal lors de son
déplacement dans le tube.

À la sortie du canon, le projectile possède :

\begin{itemize}
    \item une vitesse de translation ;
    \item une vitesse de rotation.
\end{itemize}

Ces deux mouvements coexistent pendant le vol.

\section{Le pas de rayure}

Le pas de rayure correspond à la distance nécessaire pour effectuer un tour
complet.

Par exemple :

\begin{itemize}
    \item 1 tour en 10 pouces ;
    \item 1 tour en 12 pouces ;
    \item 1 tour en 7 pouces.
\end{itemize}

Un pas plus court produit généralement une vitesse de rotation plus élevée.

\section{Calcul simplifié de la vitesse de rotation}

La vitesse de rotation dépend :

\begin{itemize}
    \item de la vitesse initiale ;
    \item du pas de rayure.
\end{itemize}

Pour une première approximation :

\begin{equation}
N=\frac{V}{P}
\end{equation}

où :

\begin{itemize}
    \item $N$ représente le nombre de tours par seconde ;
    \item $V$ représente la vitesse du projectile ;
    \item $P$ représente le pas de rayure.
\end{itemize}

\section{Ordres de grandeur}

Les vitesses de rotation obtenues sont considérables.

Un projectile moderne peut atteindre :

\begin{itemize}
    \item 100\,000 rpm ;
    \item 200\,000 rpm ;
    \item 300\,000 rpm ;
\end{itemize}

voire davantage.

Ces valeurs dépassent largement celles rencontrées dans la plupart des
systèmes mécaniques usuels.

\section{Le gyroscope}

Un gyroscope est un corps en rotation rapide.

Il possède une propriété remarquable :

\begin{quote}
Il résiste aux changements de son orientation.
\end{quote}

Cette propriété résulte de la conservation du moment cinétique.

\section{Exemple de la toupie}

Une toupie immobile tombe rapidement.

Une toupie en rotation reste verticale beaucoup plus longtemps.

La rotation modifie profondément son comportement dynamique.

Le projectile exploite le même principe physique.

\section{Moment cinétique}

Un projectile en rotation possède un moment cinétique.

Cette grandeur dépend notamment :

\begin{itemize}
    \item du moment d'inertie ;
    \item de la vitesse de rotation.
\end{itemize}

Plus le moment cinétique est important, plus il devient difficile de modifier
l'orientation du projectile.

\section{Action des forces aérodynamiques}

Pendant le vol, l'écoulement d'air exerce des forces sur le projectile.

Ces forces tendent à modifier son orientation.

Sans rotation, certaines configurations deviendraient rapidement instables.

La rotation limite fortement cette tendance.

\section{Stabilisation gyroscopique}

La stabilisation gyroscopique résulte de la combinaison :

\begin{itemize}
    \item de la rotation ;
    \item du moment cinétique ;
    \item des forces aérodynamiques.
\end{itemize}

Le projectile conserve ainsi une orientation proche de celle imposée lors de
sa sortie du canon.

\section{Pourquoi la balle ne pointe-t-elle pas exactement sa trajectoire ?}

Un projectile réel ne reste jamais parfaitement aligné avec sa trajectoire.

De faibles écarts apparaissent continuellement.

La stabilisation gyroscopique limite ces écarts mais ne les supprime pas
complètement.

Ces phénomènes conduisent notamment à :

\begin{itemize}
    \item la précession ;
    \item la nutation.
\end{itemize}

\section{Sous-stabilisation}

Si la rotation est insuffisante :

\begin{itemize}
    \item le projectile devient instable ;
    \item l'incidence augmente ;
    \item la précision se dégrade.
\end{itemize}

Dans les cas extrêmes, le projectile peut basculer complètement.

On parle alors parfois de culbute.

\section{Surstabilisation}

Une rotation excessive n'est pas toujours idéale.

Elle peut entraîner :

\begin{itemize}
    \item une augmentation de certaines contraintes mécaniques ;
    \item certains comportements aérodynamiques indésirables ;
    \item une sensibilité accrue à certains défauts.
\end{itemize}

Cependant, les effets d'une légère surstabilisation sont souvent moins
critiques que ceux d'une sous-stabilisation.

\section{Facteurs influençant la stabilité}

La stabilité gyroscopique dépend notamment :

\begin{itemize}
    \item du calibre ;
    \item de la longueur du projectile ;
    \item de sa masse ;
    \item du pas de rayure ;
    \item de la vitesse initiale ;
    \item des conditions atmosphériques.
\end{itemize}

Ces paramètres sont fortement liés entre eux.

\section{Facteur de stabilité}

Les ingénieurs utilisent souvent un facteur de stabilité gyroscopique.

Ce coefficient permet d'évaluer si la rotation est suffisante.

Plusieurs méthodes de calcul existent.

Nous les étudierons dans un chapitre ultérieur consacré aux critères de
stabilité.

\section{Rotation et traînée}

La rotation influence relativement peu la traînée principale du projectile.

Son rôle principal consiste à assurer la stabilité.

Les gains aérodynamiques proviennent surtout de la géométrie.

\section{Rotation et simulation}

Les modèles de masse ponctuelle ignorent généralement la rotation.

Les modèles avancés la prennent explicitement en compte.

Elle devient indispensable pour reproduire :

\begin{itemize}
    \item la précession ;
    \item la nutation ;
    \item le spin drift ;
    \item certains effets gyroscopiques.
\end{itemize}

\begin{simulation}

Dans les modèles 6DOF, la rotation du projectile constitue une variable
d'état fondamentale.

\end{simulation}

\section{Vers la précession}

La stabilisation gyroscopique explique pourquoi le projectile ne se renverse
pas.

Cependant, un projectile en rotation ne réagit pas aux forces comme un objet
ordinaire.

Sa réponse présente des comportements parfois surprenants.

Le premier d'entre eux est la précession.

\section{À retenir}

\begin{aretenir}

\begin{itemize}
    \item Les rayures du canon imposent une rotation au projectile.
    \item Cette rotation crée un effet gyroscopique stabilisateur.
    \item Le moment cinétique tend à conserver l'orientation du projectile.
    \item Une rotation insuffisante entraîne une instabilité.
    \item La stabilisation gyroscopique constitue le principal mécanisme de
          stabilité des balles modernes.
\end{itemize}

\end{aretenir}
